\section{今井倫太研・小野寺佳成さんの発表について}

\subsection{多人数会話における関係性ネットワークの推定とロボットによる介入戦略の提案}
本発表では，LLMを用いて多人数の自由会話から関係性ネットワークを推定し，心理的安定性の高いコミュニティ形成を目指してロボットが介入する手法（MultiRelBot）が提案された．バランス理論を活用し，疎外者のいないグループの実現が目的である．

\subsubsection{質問の背景と意図}
提案されたロボット介入には言語的発話が用いられていたが，視線や身体動作といった非言語的手段との違いや，その選択理由をより明確に理解したいと考えた．また，発表では「心理的安定性の高いコミュニティ」を目指すとあったが，これがどのように定量的・客観的に評価されているのかを確認したかった．

\subsubsection{質問内容}
\begin{itemize}
	\item 本研究においてロボットが発話という手段を用いる意義は何か？
	\item 心理的安定性の評価において，客観的な定義や測定方法は確立されているか？
\end{itemize}

\subsubsection{回答内容}
\begin{itemize}
	\item 発話による言語的介入は，関係性や文脈に対して柔軟かつ間接的な橋渡しを行える点で有効であり，非言語的な場の調整とは異なる新規性があるとのことだった．
	\item 心理的安定性の客観的定義として，ハイダーのバランス理論が採用されており，三者関係の符号の組み合わせに基づきバランス状態を定義できる．さらに，安定構造の割合を用いて定量的に評価する枠組みも整っているとの説明だった．
\end{itemize}

\subsubsection{考察}
本研究は言語的介入によって関係性そのものに働きかける点に新規性があり，ロボットによる関係調整の可能性を拡張していると感じた．心理的安定性の定量的定義も明確であり，今後の評価や応用にも期待できる．非言語的手段との比較や組み合わせについても今後検討の余地があると考えた．